% Generated draft: requires mathematical and source review.
\documentclass[11pt]{article}
\usepackage[margin=25mm]{geometry}
\usepackage{fontspec}
\setmainfont{Latin Modern Roman}
\usepackage{amsmath,amssymb,mathtools,unicode-math}
\setmathfont{Latin Modern Math}
\usepackage{enumitem,needspace,xurl,hyperref}
\setlength{\parindent}{0pt}
\setlength{\parskip}{5pt}
\setlength{\emergencystretch}{3em}
\newcommand{\N}{\mathbb{N}}
\newcommand{\Z}{\mathbb{Z}}
\newcommand{\Q}{\mathbb{Q}}
\newcommand{\R}{\mathbb{R}}
\newcommand{\C}{\mathbb{C}}
\newcommand{\F}{\mathbb{F}}
\newcommand{\A}{\mathbb{A}}
\newcommand{\PP}{\mathbb{P}}
\newcommand{\E}{\mathbb{E}}
\newcommand{\eps}{\varepsilon}
\newcommand{\dd}{\,\mathrm{d}}
\begin{document}
\section*{UnsolvedMath ID: 20003316}
\textbf{Source problem code:} catalogue entry
\subsection{Definitions and mathematical statement}
This is the first half of a continuation sentence about Conjecture 33, not an independent conjecture. It literally prints ``Conjectures 32'', a grammatical error reported as present in the source. Companion records recover the antecedent and following explanation.

Let $X_3=\mathrm{SL}_3(\R)/\mathrm{SL}_3(\Z)$, let $A$ be the full determinant-one diagonal subgroup, and put
\[
 V=\{h=(h_1,h_2,h_3)\in\R^3:h_1+h_2+h_3=0\},\qquad
 R_h=\{\operatorname{diag}(e^{th_1},e^{th_2},e^{th_3}):t\ge0\}
 \quad(h\in V\setminus\{0\}).
\]
Conjecture 33 asks, for every $h,k\ne0$, for $x\in X_3$ such that both $R_hx$ and $R_kx$ are relatively compact but $Ax$ is not relatively compact. The historical same-weight-cone argument applies when both rays belong to
\[
 \{\operatorname{diag}(e^u,e^v,e^{-u-v}):u,v\ge0\},
\]
with analogous coordinate permutations and simultaneous duality; it uses the weighted intersection assertion of Conjecture 32. The source then attributes to D. Kleinbock an opposite-Weyl-chamber argument, continued in the next record, giving witnesses and full Hausdorff dimension. Its exact historical reference is [108], not a general two-ray theorem asserted here.

The aid reports a useful geometric distinction. The coordinate walls $h_i=0$ define six weight cones in $V$, while the root walls $h_i=h_j$ define six Weyl chambers. These two hexagonal fans are offset by $30$ degrees and have a twelve-sector common refinement. For regular generators, meaning $h_i\ne h_j$ for $i\ne j$, relative chamber positions have four possible cyclic distances $0,1,2,3$ after common coordinate permutation, simultaneous lattice duality $h\mapsto-h$, and ray exchange. This classifies chamber-position types, not the continuously varying rays themselves. Common weight-cone pairs can only have Weyl distance $0$ or $1$. Opposite chambers have distance $3$, equivalently
\[
 (h_i-h_j)(k_i-k_j)<0\quad\text{for every }i<j.
\]
The pair $h=(2,1,-3)$, $k=(3,-1,-2)$ has the same coordinate ordering but does not lie in a common weight cone, even allowing the stated common symmetries; it is also not opposite. Thus the two historical routes are not exhaustive. Coordinate-wall endpoints $h_i=0$ and singular root walls $h_i=h_j$ are distinct boundary phenomena and should not be conflated.

The attached aid reports a modern resolution for arbitrary finite or countable families of nontrivial diagonal rays. The joint-bounded-orbit set from An--Guan--Kleinbock is thick in the eight-dimensional $X_3$, whereas the full-$A$-bounded locus has Hausdorff dimension two by Einsiedler--Katok--Lindenstrauss. Removing that lower-dimensional locus leaves a thick set, of dimension eight in each nonempty open set, on which all prescribed ray orbits are bounded and the full diagonal orbit is unbounded. These are supplied literature claims requiring independent verification, distinct from the historical special cases.
\subsection{Short English statement}
This continuation explains two historical routes to bounded orbits along two rays with an unbounded full diagonal orbit. Weight cones and Weyl chambers differ, so those routes do not cover every pair; the attached aid reports a later countable-ray full-dimension resolution.
\subsection{Sources}
\begin{itemize}
\item American Institute of Mathematics, \emph{Emerging applications of measure rigidity}, continuation of Conjecture 33, with D. Kleinbock's opposite-chamber observation and [108], \url{https://aimath.org/WWN/measrigid/measrigid.pdf}. The supplied companion summary identifies [108] with Kleinbock--Margulis but does not include the full citation.
\item Companion local records in \texttt{aim-other-notes.json}: ID $20003315$, \texttt{AIM-OTHER-0122}, \texttt{notes[121]}, supplies Conjecture 33 and the closed weight cone; ID $20003317$, \texttt{AIM-OTHER-0124}, \texttt{notes[123]}, supplies the opposite-chamber continuation; ID $20003313$, \texttt{AIM-OTHER-0120}, \texttt{notes[119]}, states the weighted intersection Conjecture 32.
\item The supplied aid attributes modern joint-bounded thickness to An--Guan--Kleinbock and dimension two of the full-diagonal bounded locus to Einsiedler--Katok--Lindenstrauss; full bibliographic references are absent.
\item Local UnsolvedMath registry, permanent ID $20003316$, source code \texttt{catalogue entry}, canonical \texttt{aim-other-notes.json}, \texttt{notes[122]}, dated 14 August 2026. Research attempt 1: context only, ``not a problem'' at run, proved rigor, low novelty confidence; machine-generated and not peer reviewed.
\end{itemize}
\end{document}
